\chapter{When the yardstick changes}\label{ch:moving-field}

\section{Learning about the world is not the same as changing it}
A community can improve a water reserve by moving water. It can also learn that its previous reserve target was unsuitable. The first changes the physical state; the second changes the field used to interpret that state. Both can be important, but only one is the physical action described by the original receipt.

Let $V(x,\theta)$ be a family of potentials, with $\theta$ identifying references, scales and other declared field parameters. A fixed-field action still has $E_\theta=V(x_0,\theta)-V(x_1,\theta)$. Once $\theta$ changes, an account spanning both versions needs to show that change explicitly.

\section{The full differential has two components}
Assume $V$ is $C^1$ along a regular path $(x(t),\theta(t))$. The chain rule gives
\begin{equation}\label{eq:moving-differential}
 dV=\nabla_xV\cdot dx+\partial_\theta V\cdot d\theta.
\end{equation}
Consequently,
\begin{equation}\label{eq:moving-integral}
 V(x_1,\theta_1)-V(x_0,\theta_0)
 =\int\nabla_xV\cdot dx+\int\partial_\theta V\cdot d\theta.
\end{equation}
The total is an exact state-function difference on the enlarged state-and-parameter space. The two component integrals are generally path dependent separately. One cannot discard the parameter component and retain a universal endpoint theorem for the physical-action component across changing fields.

For example, with $V(x,\theta)=\theta x$, the physical component is $\theta\,dx$ and the parameter component is $x\,d\theta$. Moving $x$ first and changing $\theta$ second allocates different amounts to the two components than doing them in reverse order. Both sums give the same final minus initial product.

\section{A finite version makes the ordering explicit}
For a transition from $(x,\theta)$ to $(y,\theta')$, insert the intermediate value $V(y,\theta)$:
\begin{align}\label{eq:field-split}
 V(x,\theta)-V(y,\theta')
 ={}&[V(x,\theta)-V(y,\theta)]\nonumber\\
   &+[V(y,\theta)-V(y,\theta')].
\end{align}
The first bracket values the physical change under the original field. The second records revaluation at the final physical state. An alternative ordering evaluates the physical change under the new field and revalues at the initial state. Both are valid decompositions with different meanings. A protocol must choose and record its ordering.

At $(14,6)$, changing both Gaussian scales from two to four changes potential from four to one without moving water. That three-unit difference is a field-change entry. An actor who was idle during the revision has not performed a three-unit restoration.

\section{A closed physical loop can exploit an untracked scale change}
Let $W(0)=0$, $W(1)=1$, and define fields $V_A=2W$, $V_B=W$. An action $1\to0$ under A gives $+2$. A later reverse action $0\to1$ under B gives $-1$. Returning the field to A completes a loop in both physical state and field version, but the two action entries alone sum to $+1$.

The missing field-change entries reconcile the loop. No material was created and the fixed-field theorem was not violated: its premise of one field did not hold for the action-only sum. Treating those heterogeneous entries as a persistent spendable claim without a bridge would create an accounting opportunity from changes in the ruler.

This example identifies a concrete institutional risk without deciding a universal remedy. A canonical capacity reference, an explicit conversion rule or separated accounts may each be considered in a future design. Each needs a justified meaning and its own analysis. The book does not adopt one by naming the problem.

\section{Historical receipts retain their original meaning}
A later field can provide a useful retrospective diagnostic, but it should not overwrite what the original account said. Keep the original state, field version and signed value; append the new calculation with its purpose and authority. Otherwise genuine physical improvement and a change in interpretation become impossible to distinguish.

This matters for trust. A community should be able to learn that a reference was wrong without being told that its historical work has silently changed units. An institution should also be able to correct a mistaken measurement without erasing the earlier evidence. Measurement correction, physical reversal and field revision are separate events.

\section{A common coefficient across fields remains a hypothesis}
The canonical benchmark gave $J_{\mathrm{rel}}=V$ under a particular energy normalization. For more general fields, one might investigate
\begin{equation}
 J_{{\mathrm{rel}},\theta}(x)\overset{?}{=}
 \beta_{\mathrm{bridge}}V_\theta(x)
\end{equation}
with one coefficient across the specified fields. The question is stronger than agreement within one fitted field. It requires independently defined potentials, compatible density measures and a justified common denomination.

A curvature comparison $K_\theta\overset{?}{=}\beta_{\mathrm{bridge}}H_\theta$ can be part of that programme, with accessible tangent directions treated correctly. Local curvature agreement alone does not establish a global identity or a universal conversion among resources. Nor may the potential branch be defined from the probability branch and then presented as independent evidence of agreement.

\section{Why a fixed yardstick is not the final social answer}
Freezing a field forever would simplify accounting, but knowledge and physical conditions change. Drought can alter feasible reserve levels; technology can change the relationship between energy and service; new evidence can expose an omitted harm. A trustworthy system needs a way to improve its model without presenting every revision as an actor's achievement.

The positive opportunity is an account that remains both relevant and inspectable. Versioned fields can make learning visible, while separated revaluation entries preserve the meaning of physical history. That could support more informed public decisions and prevent windfalls caused solely by hidden scale changes. The rules governing persistent claims remain an institutional research problem.

The next chapter turns from a changing field to imperfect knowledge of a fixed one. Even when the declaration is held constant, measurements and shared calibration can make an interaction uncertain.
